\chapter{SX-WSA1 and SX-WSA1R}\label{ch:models}

The SX-WSA1 keyboard and the SX-WSA1R rack module run the same program and speak the same
System Exclusive protocol. Two features are nonetheless available on the keyboard and not
on the rack, in both directions.

\begin{center}
\begin{tabular}{lcc}
\toprule
 & SX-WSA1 & SX-WSA1R \\
 & (keyboard) & (rack) \\
\midrule
message framing, checksum, handshake & yes & yes \\
\textsc{sound} bulk dump & yes & yes \\
\textsc{combination} bulk dump & yes & yes \\
\textsc{system, part \& midi} bulk dump & yes & yes \\
dump requests & yes & yes \\
General MIDI On / Off & yes & yes \\
\midrule
\textsc{sequencer} bulk dump & yes & \textbf{no} \\
tempo message (family \bytes{25}) & yes & \textbf{no} \\
\bottomrule
\end{tabular}
\end{center}

\begin{caution}{On the rack, a SEQUENCER dump sends nothing and says nothing}
The \textsc{sequencer} row is present on the rack's \textsc{sysex bulk dump} menu, and
selecting it and pressing \textsc{send} appears to work. Nothing is transmitted --- not
even the end-of-category message that closes every other transfer. A librarian waiting for
data will simply time out.

The same applies to the \textsc{sequencer} portion of \textsc{total keyboard}.
\end{caution}

The tempo message is refused in both directions on the rack: it will neither act on one
received nor send one.

\begin{note}{The model triple does not tell you which you are talking to}
Both models transmit the same model triple, so it cannot be used to distinguish them. The
only reliable test over MIDI is behavioural --- request a \textsc{sequencer} dump and see
whether anything arrives.
\end{note}
